\begin{center}
\LARGE
Global projection operators for the Cylindrical Algebraic
                   Decomposition algorithm

\end{center}

\begin{center}
\Large
Emma-Neila Gonzalez-Campos, 
                      Laureano Gonzalez-Vega

\end{center}

\noindent
Date:  July 19th (Friday)\\
Time:  14:40-15:00 \\

\begin{center}
\large
Abstract
\end{center}

The projection phase in the CAD algorithm proceeds by eliminating one variable on a given set of
polynomials X via the computation of the principal subresultant coefficients of a well precised set
of polynomials pairs on X (including their derivatives and their reductums). Since the size, at least
in number, of this new set of polynomials is usually very big, any improvement in the projection
phase would convey to dramatically speed up the efficiency of the CAD algorithm. The purpose of
this talk is to present two approaches allowing, in some cases, to simplify the projection phase in
the CAD algorithm: 
\begin{enumerate}
\item The first one tries to avoid the consideration of all the different pairs of polynomials and their
reductums by parametrizing the rank of a well precised matrix with polynomial entries. 
\item The second one tries to perform an improved projection phase through the elimination of blocks
of variables instead of the usual one by one variable elimination. 
\end{enumerate}
The theoretical tools to be used include Barnett's Theorem characterizing the degree of the gcd of
a finite set of polynomials through the rank of a well-precised matrix and the parametrization of
Hermite's Method for counting the number of real solutions of a polynomial system of equations. 

